\section{Quotient frames and involutions}

The three quotient frames are the orbit partitions of the three
involutions

\[
\tau_0,
\qquad
\tau_1,
\qquad
\tau_2
\]

in the complementary group

\[
C\cong S_3.
\]

Their structural data are summarized in
Table~\ref{tab:frame-summary}.

\begin{table}[h]
\centering
\begin{tabular}{llll}
\toprule
Frame & Block swap & Quotient & Stabilizer\\
\midrule
\(\mathcal P_0\) & \(\tau_0\) & \(L(P)\) & \(S_5\times C_2\)\\
\(\mathcal P_1\) & \(\tau_1\) & \(L(P)\) & \(S_5\times C_2\)\\
\(\mathcal P_2\) & \(\tau_2\) & \(L(P)\) & \(S_5\times C_2\)\\
\bottomrule
\end{tabular}
\caption{The certified quotient-frame orbit.}
\label{tab:frame-summary}
\end{table}

The induced action on the three frames is

\[
\rho(\tau_0)=[0,2,1],
\]

\[
\rho(\tau_1)=[2,1,0],
\]

and

\[
\rho(\tau_2)=[1,0,2].
\]

The frames are pairwise block-disjoint and therefore contain forty-five
distinct unordered pairs.

\subsection{The frame \(\mathcal P_0\)}

The certified block partition is

\[
\mathcal P_0
=
\left\{
  \{0,1\}, \{2,3\}, \{4,5\}, \{6,7\}, \{8,9\},
  \{10,11\}, \{12,13\}, \{14,15\}, \{16,17\}, \{18,19\},
  \{20,21\}, \{22,23\}, \{24,25\}, \{26,27\}, \{28,29\}
\right\}.
\]

Its block-swap involution is encoded in image-list notation by

\[
\tau_0
=
\left[
  1, 0, 3, 2, 5, 4, 7, 6, 9, 8,
  11, 10, 13, 12, 15, 14, 17, 16, 19, 18,
  21, 20, 23, 22, 25, 24, 27, 26, 29, 28
\right].
\]

The certificate verifies

\[
\tau_0^2=1,
\qquad
\operatorname{Orb}(\langle\tau_0\rangle)
=
\mathcal P_0,
\]

and records the induced frame action

\[
\rho(\tau_0)
=
[0, 2, 1].
\]

The quotient has fifteen vertices and thirty edges, and the covering-edge
test passes for this partition.

\subsection{The frame \(\mathcal P_1\)}

The certified block partition is

\[
\mathcal P_1
=
\left\{
  \{0,11\}, \{1,28\}, \{2,25\}, \{3,12\}, \{4,19\},
  \{5,26\}, \{6,23\}, \{7,14\}, \{8,17\}, \{9,20\},
  \{10,29\}, \{13,24\}, \{15,22\}, \{16,21\}, \{18,27\}
\right\}.
\]

Its block-swap involution is encoded in image-list notation by

\[
\tau_1
=
\left[
  11, 28, 25, 12, 19, 26, 23, 14, 17, 20,
  29, 0, 3, 24, 7, 22, 21, 8, 27, 4,
  9, 16, 15, 6, 13, 2, 5, 18, 1, 10
\right].
\]

The certificate verifies

\[
\tau_1^2=1,
\qquad
\operatorname{Orb}(\langle\tau_1\rangle)
=
\mathcal P_1,
\]

and records the induced frame action

\[
\rho(\tau_1)
=
[2, 1, 0].
\]

The quotient has fifteen vertices and thirty edges, and the covering-edge
test passes for this partition.

\subsection{The frame \(\mathcal P_2\)}

The certified block partition is

\[
\mathcal P_2
=
\left\{
  \{0,29\}, \{1,10\}, \{2,13\}, \{3,24\}, \{4,27\},
  \{5,18\}, \{6,15\}, \{7,22\}, \{8,21\}, \{9,16\},
  \{11,28\}, \{12,25\}, \{14,23\}, \{17,20\}, \{19,26\}
\right\}.
\]

Its block-swap involution is encoded in image-list notation by

\[
\tau_2
=
\left[
  29, 10, 13, 24, 27, 18, 15, 22, 21, 16,
  1, 28, 25, 2, 23, 6, 9, 20, 5, 26,
  17, 8, 7, 14, 3, 12, 19, 4, 11, 0
\right].
\]

The certificate verifies

\[
\tau_2^2=1,
\qquad
\operatorname{Orb}(\langle\tau_2\rangle)
=
\mathcal P_2,
\]

and records the induced frame action

\[
\rho(\tau_2)
=
[1, 0, 2].
\]

The quotient has fifteen vertices and thirty edges, and the covering-edge
test passes for this partition.

\subsection{Certificate checks and boundary}

The certificate records

\[
|\operatorname{Aut}(X)|=720,
\qquad
|\operatorname{im}\rho|=6,
\qquad
|\ker\rho|=120,
\]

and verifies that the three partitions contain fifteen blocks each, that
all three block swaps are fixed-point-free involutions, and that every
orbit partition agrees with its displayed block list.

The three block sets are pairwise disjoint. Hence

\[
\left|
\mathcal P_0
\cup
\mathcal P_1
\cup
\mathcal P_2
\right|
=
45.
\]

The machine-readable source is

\begin{verbatim}
sources/project41-paper42/project42_quotient_frame_orbit_certificate_004.json
\end{verbatim}

Its SHA-256 digest is

\begin{verbatim}
9e6f11edfffa8cb1fbf4db86935d9e28c0442a1bd72fbe21607055423b66d041
\end{verbatim}

This certificate proves the three-member automorphism orbit of the
natural quotient frame. It does not claim that every admissible regular
\(L(P)\)-quotient frame belongs to this orbit. The source boundary states:

\begin{quote}
No exhaustive census of admissible quotient systems outside the natural automorphism orbit is claimed.
\end{quote}

